% Compilar dos veces con pdflatex.
\documentclass[aspectratio=169,11pt]{beamer}
\usepackage[utf8]{inputenc}
\usepackage[T1]{fontenc}
\usepackage[spanish,es-nodecimaldot]{babel}
\usepackage{amsmath,amssymb,booktabs,tabularx,array,tikz}
\usetikzlibrary{arrows.meta,positioning,shapes.geometric,fit}

\definecolor{UAPNavy}{HTML}{17324D}
\definecolor{UAPTeal}{HTML}{007F82}
\definecolor{UAPGold}{HTML}{E2A33A}
\definecolor{UAPLight}{HTML}{F3F6F8}
\definecolor{UAPGray}{HTML}{536471}

\tikzset{
  flow/.style={-{Stealth[length=2.4mm]},thick,draw=UAPNavy},
  state/.style={circle,draw=UAPTeal,fill=UAPTeal!12,very thick,minimum size=9mm,font=\bfseries},
  goal/.style={state,double,double distance=1.2pt,draw=UAPGold},
  card/.style={draw=UAPTeal,fill=UAPLight,rounded corners=3pt,align=center,minimum height=11mm},
  data/.style={draw=UAPGold,fill=UAPGold!15,rounded corners=2pt,align=center,minimum height=9mm},
  searchnode/.style={draw=UAPNavy,fill=white,rounded corners=2pt,align=center,minimum width=12mm,minimum height=8mm}
}

\usetheme{Boadilla}
\setbeamercolor{normal text}{fg=UAPNavy,bg=white}
\setbeamercolor{structure}{fg=UAPTeal}
\setbeamercolor{frametitle}{fg=white,bg=UAPNavy}
\setbeamercolor{title}{fg=white}
\setbeamercolor{subtitle}{fg=UAPGold}
\setbeamercolor{block title}{fg=white,bg=UAPTeal}
\setbeamercolor{block body}{fg=UAPNavy,bg=UAPLight}
\setbeamercolor{alerted text}{fg=UAPTeal}
\setbeamerfont{frametitle}{series=\bfseries}
\setbeamertemplate{navigation symbols}{}
\setbeamertemplate{itemize item}{\color{UAPGold}$\blacktriangleright$}
\setbeamertemplate{itemize subitem}{\color{UAPTeal}$\bullet$}
\setbeamertemplate{footline}{%
  \leavevmode\hbox{%
    \begin{beamercolorbox}[wd=.78\paperwidth,ht=2.8ex,dp=1.2ex,leftskip=1em]{author in head/foot}%
      \color{white}\insertshorttitle
    \end{beamercolorbox}%
    \begin{beamercolorbox}[wd=.22\paperwidth,ht=2.8ex,dp=1.2ex,rightskip=1em plus1fil]{date in head/foot}%
      \color{UAPGray}\hfill Semana 3 \enspace | \enspace \insertframenumber/\inserttotalframenumber
    \end{beamercolorbox}}}

\newcommand{\concepto}[1]{\textcolor{UAPTeal}{\textbf{#1}}}
\newcommand{\pregunta}[1]{\begin{block}{Pregunta clave}#1\end{block}}
\title[IA · Clase 3]{Búsqueda informada y heurísticas}
\subtitle{De costo uniforme a A*: diseño, garantías y evaluación}
\author{Curso de Inteligencia Artificial}
\institute{Semana 3 · Clase 3}
\date{}

\begin{document}

% 1
\begin{frame}[plain]
  \begin{beamercolorbox}[wd=\paperwidth,ht=\paperheight,sep=2em]{frametitle}
    \vfill {\Huge\bfseries\inserttitle\par}\vspace{.5em}
    {\color{UAPGold}\Large\insertsubtitle\par}\vspace{2em}
    {\color{white}\large\insertauthor\par\insertinstitute\par}\vfill
  \end{beamercolorbox}
\end{frame}

% 2
\begin{frame}{Propósito de la clase}
  Incorporar conocimiento del problema para orientar la búsqueda sin confundir rapidez con garantía.
  \begin{block}{Al finalizar, podremos}
    \begin{itemize}
      \item definir y auditar una función heurística;
      \item ejecutar búsqueda voraz y A*;
      \item distinguir admisibilidad de consistencia;
      \item explicar cuándo A* requiere reaperturas;
      \item diseñar una comparación reproducible entre UCS y A*.
    \end{itemize}
  \end{block}
\end{frame}

% 3
\begin{frame}{Punto de partida}
  La clase anterior dejó establecidos:
  \begin{itemize}
    \item el problema $P=(S,A,T,s_0,G,c)$;
    \item nodo, frontera y conjunto explorado;
    \item control de ciclos y mejores costos;
    \item BFS, DFS y costo uniforme;
    \item garantías condicionadas por representación y costos.
  \end{itemize}
  \begin{block}{Continuidad}
    UCS ordena por costo recorrido; ahora agregaremos una estimación del costo restante.
  \end{block}
\end{frame}

% 4
\begin{frame}{Pregunta de apertura}
  \pregunta{Si dos caminos cuestan lo mismo hasta ahora, ¿cómo decidir cuál explorar primero?}
  \begin{itemize}
    \item ¿Qué conocimiento adicional está permitido?
    \item ¿En qué unidades debe expresarse?
    \item ¿Cómo sabemos si orienta sin engañar?
    \item ¿Qué garantía queremos conservar?
  \end{itemize}
\end{frame}

% 5
\begin{frame}[fragile]{El esquema de búsqueda no cambia}
\small
\begin{verbatim}
insertar nodo inicial en frontera
mientras frontera no esté vacía:
    extraer el nodo con menor valor de prioridad
    descartar entradas obsoletas
    si es objetivo: devolver camino
    relajar cada transición legal
devolver fracaso
\end{verbatim}
  \alert{La estrategia cambia al definir la prioridad de la frontera.}
\end{frame}

% 6
\begin{frame}{Una familia, distintas prioridades}
  \small
  \begin{tabularx}{\textwidth}{>{\bfseries}p{.18\textwidth}p{.20\textwidth}X}
    \toprule
    Estrategia & Prioridad & Pregunta que privilegia \\\midrule
    BFS & profundidad & ¿Cuántos pasos llevo? \\
    UCS & $g(n)$ & ¿Cuánto costó llegar? \\
    Voraz & $h(n)$ & ¿Qué parece más cerca? \\
    A* & $g(n)+h(n)$ & ¿Qué solución total parece más barata? \\
    \bottomrule
  \end{tabularx}
  \vfill
  La prioridad no modifica el grafo ni el costo real del camino.
\end{frame}

% 7
\begin{frame}{Búsqueda no informada e informada}
  \small
  Ambas parten de $P=(S,A,T,s_0,G,c)$. La diferencia es la información usada para ordenar la frontera.
  \vspace{.4em}
  \begin{columns}[T]
    \column{.48\textwidth}
      \begin{block}{No informada: sin $h$}
        \begin{itemize}
          \item Usa la definición del problema y el costo recorrido.
          \item No estima el costo que falta.
          \item Incluye BFS, DFS, limitada, iterativa y UCS.
          \item Prioriza por profundidad, orden o $g(n)$.
        \end{itemize}
      \end{block}
      {\scriptsize UCS pertenece aquí: $g(n)$ es costo exacto del pasado, no una estimación del futuro.}
    \column{.48\textwidth}
      \begin{block}{Informada: incorpora $h$}
        \begin{itemize}
          \item Añade una estimación dirigida al objetivo.
          \item Voraz prioriza $h(n)$.
          \item A* prioriza $g(n)+h(n)$.
          \item La garantía depende de la estrategia y de $h$.
        \end{itemize}
      \end{block}
      {\scriptsize Voraz es informada, pero no es óptima en general.}
  \end{columns}
  \vfill
  \alert{No informada no significa aleatoria; informada no significa conocer el futuro ni garantizar optimalidad.}
\end{frame}

% 8
\begin{frame}{Punto de control: costo uniforme}
  UCS extrae el nodo con menor $g(n)$.
  \begin{itemize}
    \item Usa una cola de prioridad.
    \item Conserva el mejor costo conocido por estado.
    \item Prueba el objetivo al extraer, no al generar.
    \item Descarta entradas obsoletas.
    \item Es óptima bajo los supuestos declarados.
  \end{itemize}
  \begin{alertblock}{Límite}
    Puede expandir muchas alternativas baratas que se alejan del objetivo.
  \end{alertblock}
\end{frame}

% 9
\begin{frame}{¿Qué añade la búsqueda informada?}
  Añade conocimiento para ordenar la exploración:
  \[
    h(n)\approx \text{costo mínimo que falta desde }n\text{ hasta un objetivo}
  \]
  \begin{itemize}
    \item No altera $g(n)$.
    \item No crea conexiones inexistentes.
    \item No corrige estados insuficientes.
    \item No reemplaza restricciones.
  \end{itemize}
  \alert{Una heurística guía la búsqueda dentro del problema ya formulado.}
\end{frame}

% 10
\begin{frame}{Función heurística}
  Una heurística es una función:
  \[
    h:S\rightarrow\mathbb{R}
  \]
  \begin{itemize}
    \item $h(n)$: estimación del costo restante;
    \item $h^*(n)$: costo óptimo restante verdadero;
    \item $h(g)=0$ para todo estado objetivo $g$.
  \end{itemize}
  \begin{block}{Convención}
    La heurística se evalúa sobre estados, aunque se escriba $h(n)$ por conveniencia.
  \end{block}
\end{frame}

% 11
\begin{frame}{Tres cantidades diferentes}
  \small
  \begin{tabularx}{\textwidth}{>{\bfseries}p{.15\textwidth}Xp{.28\textwidth}}
    \toprule
    Cantidad & Significado & Disponibilidad \\\midrule
    $g(n)$ & costo acumulado desde $s_0$ & conocido al generar \\
    $h(n)$ & estimación del costo restante & calculado por la heurística \\
    $f(n)$ & prioridad de la estrategia & depende del algoritmo \\
    \bottomrule
  \end{tabularx}
  \vfill
  En A*:
  \[
    \boxed{f(n)=g(n)+h(n)}
  \]
  $f$ estima el costo total de una solución que pasa por $n$.
\end{frame}

% 12
\begin{frame}{La unidad debe coincidir}
  Si $g$ está en minutos, $h$ también debe estar en minutos.
  \begin{table}
  \begin{tabularx}{.88\textwidth}{>{\bfseries}p{.28\textwidth}X}
    \toprule
    Costo objetivo & Heurística candidata \\\midrule
    distancia & separación geométrica inferior \\
    tiempo & distancia / velocidad máxima válida \\
    cantidad de pasos & mínimo de pasos pendientes \\
    \bottomrule
  \end{tabularx}
  \end{table}
  \begin{alertblock}{Inválido}
    Sumar kilómetros y minutos sin una conversión justificada.
  \end{alertblock}
\end{frame}

% 13
\begin{frame}{Dos extremos útiles}
  \begin{columns}[T]
    \column{.47\textwidth}
      \begin{block}{Heurística nula}
        \[
          h(n)=0
        \]
        A* se convierte en UCS. Es segura, pero no aporta orientación.
      \end{block}
    \column{.47\textwidth}
      \begin{block}{Heurística perfecta}
        \[
          h(n)=h^*(n)
        \]
        Orienta idealmente, pero calcularla suele equivaler a resolver el problema.
      \end{block}
  \end{columns}
  \vfill
  \centering\concepto{Buscamos una aproximación barata, informativa y justificable.}
\end{frame}

% 14
\begin{frame}{Informativa no significa segura}
  \begin{tabularx}{\textwidth}{>{\bfseries}p{.28\textwidth}X}
    \toprule
    Dimensión & Pregunta \\\midrule
    Calidad empírica & ¿Reduce expansiones y discrimina alternativas? \\
    Garantía teórica & ¿Respeta una relación demostrable con $h^*$ y los costos? \\
    \bottomrule
  \end{tabularx}
  \vfill
  \begin{alertblock}{No alcanza con parecer correcta}
    Una alta correlación con el costo verdadero no demuestra admisibilidad.
  \end{alertblock}
\end{frame}

% 15
\begin{frame}{Lista de control inicial}
  Antes de usar $h$, declarar:
  \begin{enumerate}
    \item qué costo estima y en qué unidad;
    \item qué información utiliza;
    \item por qué vale cero en el objetivo;
    \item si puede ser negativa;
    \item cómo se comprobará consistencia;
    \item cuánto cuesta calcularla.
  \end{enumerate}
  \begin{block}{Criterio}
    Una fórmula intuitiva no es todavía una heurística defendible.
  \end{block}
\end{frame}

% 16
\begin{frame}{Búsqueda voraz por el mejor primero}
  La búsqueda voraz ordena la frontera solo por:
  \[
    \boxed{f(n)=h(n)}
  \]
  \begin{itemize}
    \item Elige el estado que parece más cercano al objetivo.
    \item Ignora cuánto costó llegar hasta allí.
    \item Puede encontrar rápidamente una solución plausible.
    \item No garantiza el menor costo.
  \end{itemize}
\end{frame}

% 17
\begin{frame}{Contraejemplo para la búsqueda voraz}
  \centering
  \begin{tikzpicture}[scale=.9,transform shape]
    \node[state] (s) at (0,0) {S};
    \node[state] (a) at (3,1.35) {A};
    \node[state] (b) at (3,-1.35) {B};
    \node[goal] (g) at (6,0) {G};
    \draw[flow] (s)--node[above,sloped]{100}(a);
    \draw[flow] (s)--node[below,sloped]{1}(b);
    \draw[flow] (a)--node[above,sloped]{1}(g);
    \draw[flow] (b)--node[below,sloped]{5}(g);
    \node[data,above=4mm of a] {$h(A)=1$};
    \node[data,below=4mm of b] {$h(B)=5$};
  \end{tikzpicture}
  \vfill
  Voraz elige $A$ y devuelve costo 101; existe $S\rightarrow B\rightarrow G$ de costo 6.
\end{frame}

% 18
\begin{frame}{¿Por qué falla la voraz?}
  \begin{table}
  \begin{tabular}{lrrr}
    \toprule
    Nodo & $g$ & $h$ & prioridad voraz \\\midrule
    A & 100 & 1 & 1 \\
    B & 1 & 5 & 5 \\
    \bottomrule
  \end{tabular}
  \end{table}
  \begin{block}{Lectura}
    La heurística describe costo restante, no costo total.
  \end{block}
  \alert{Al ignorar $g$, una promesa cercana puede ocultar un pasado extremadamente caro.}
\end{frame}

% 19
\begin{frame}{Alcance de la búsqueda voraz}
  \begin{itemize}
    \item En un grafo finito con control de repetidos, termina.
    \item En espacios infinitos no es completa en general.
    \item No es óptima, incluso si $h$ es admisible.
    \item Debe evaluarse tanto velocidad como calidad de solución.
  \end{itemize}
  \begin{block}{Uso razonable}
    Obtener una solución rápida cuando la optimalidad no es obligatoria.
  \end{block}
\end{frame}

% 20
\begin{frame}{Algoritmo A*}
  A* equilibra lo recorrido y lo que falta:
  \[
    \boxed{f(n)=g(n)+h(n)}
  \]
  \begin{itemize}
    \item $g(n)$ aporta evidencia del camino ya recorrido.
    \item $h(n)$ orienta hacia regiones prometedoras.
    \item Se extrae el nodo con menor $f$.
  \end{itemize}
  \begin{block}{Caso límite}
    Con $h=0$, A* reproduce el comportamiento de UCS.
  \end{block}
\end{frame}

% 21
\begin{frame}{Qué debe guardar A*}
  Cada nodo conserva:
  \[
    n=(\text{estado},\text{padre},\text{acción},g,h,f)
  \]
  La implementación también necesita:
  \begin{itemize}
    \item cola de prioridad por $f$;
    \item \texttt{mejor\_g[estado]};
    \item regla de desempate estable;
    \item padres para reconstruir el camino;
    \item política explícita de reapertura.
  \end{itemize}
\end{frame}

% 22
\begin{frame}[fragile]{Pseudocódigo de A*}
\small
\begin{verbatim}
frontera <- prioridad por f
mejor_g[inicial] <- 0
insertar inicial con f = h(inicial)
mientras haya frontera:
    n <- extraer mínimo
    si n.g es obsoleto: continuar
    si n es objetivo: devolver camino
    para cada sucesor:
        nuevo_g <- n.g + costo
        si nuevo_g mejora mejor_g:
            registrar e insertar con f = nuevo_g + h
devolver fracaso
\end{verbatim}
\end{frame}

% 23
\begin{frame}{Grafo didáctico para A*}
  \centering
  \begin{tikzpicture}[scale=.82,transform shape]
    \node[state] (s) at (0,0) {S};
    \node[state] (a) at (2.2,1.45) {A};
    \node[state] (b) at (2.2,-1.45) {B};
    \node[state] (c) at (4.5,1.45) {C};
    \node[state] (d) at (4.5,-1.1) {D};
    \node[goal] (g) at (7,0) {G};
    \draw[flow] (s)--node[above,sloped]{2}(a);
    \draw[flow] (s)--node[below,sloped]{2}(b);
    \draw[flow] (a)--node[above]{2}(c);
    \draw[flow] (a)--node[above,sloped]{5}(d);
    \draw[flow] (b)--node[below,sloped]{2}(d);
    \draw[flow] (c)--node[above,sloped]{3}(g);
    \draw[flow] (d)--node[below,sloped]{6}(g);
    \node[font=\scriptsize,below=2mm of s] {$h=7$};
    \node[font=\scriptsize,above=2mm of a] {$h=5$};
    \node[font=\scriptsize,below=2mm of b] {$h=7$};
    \node[font=\scriptsize,above=2mm of c] {$h=3$};
    \node[font=\scriptsize,below=2mm of d] {$h=6$};
    \node[font=\scriptsize,below=2mm of g] {$h=0$};
  \end{tikzpicture}
  \vfill
  Aristas dirigidas · costos no negativos · desempate por orden de inserción.
\end{frame}

% 24
\begin{frame}{A*: inicialización}
  \begin{table}
  \begin{tabular}{lrrr}
    \toprule
    Nodo & $g$ & $h$ & $f$ \\\midrule
    S & 0 & 7 & 7 \\
    \bottomrule
  \end{tabular}
  \end{table}
  \begin{block}{Primera extracción}
    Se extrae $S$ porque es el único nodo disponible.
  \end{block}
  \centering Regla de desempate del ejemplo: \concepto{orden de inserción}.
\end{frame}

% 25
\begin{frame}{A*: después de expandir S}
  \begin{table}
  \begin{tabular}{lrrr}
    \toprule
    Frontera & $g$ & $h$ & $f$ \\\midrule
    A & 2 & 5 & 7 \\
    B & 2 & 7 & 9 \\
    \bottomrule
  \end{tabular}
  \end{table}
  \begin{block}{Decisión}
    A* extrae $A$. Ambos caminos recorridos cuestan lo mismo; la heurística rompe la igualdad.
  \end{block}
\end{frame}

% 26
\begin{frame}{A*: de A al objetivo}
  \small
  \begin{tabular}{lrrr}
    \toprule
    Frontera tras expandir A & $g$ & $h$ & $f$ \\\midrule
    C & 4 & 3 & 7 \\
    B & 2 & 7 & 9 \\
    D & 7 & 6 & 13 \\
    \bottomrule
  \end{tabular}
  \vfill
  Al expandir $C$ aparece $G$ con $g=f=7$. Al extraer $G$:
  \[
    \boxed{S\rightarrow A\rightarrow C\rightarrow G,\qquad C^*=7}
  \]
\end{frame}

% 27
\begin{frame}{UCS y A* sobre la misma instancia}
  \small
  \begin{tabularx}{\textwidth}{>{\bfseries}p{.34\textwidth}XX}
    \toprule
    Resultado & UCS & A* \\\midrule
    Camino & $S-A-C-G$ & $S-A-C-G$ \\
    Costo & 7 & 7 \\
    Prioridad & $g$ & $g+h$ \\
    Expandidos antes de extraer G & $S,A,B,C,D$ & $S,A,C$ \\
    \bottomrule
  \end{tabularx}
  \vfill
  \alert{La heurística reduce expansiones en este ejemplo; no reduce el costo óptimo.}
\end{frame}

% 28
\begin{frame}{A* no significa óptimo automáticamente}
  \small
  La garantía depende de:
  \begin{itemize}
    \item ramificación finita;
    \item costos no negativos y, para terminación, $c\geq\varepsilon>0$;
    \item objetivo comprobado al extraer;
    \item conservación del mejor $g$;
    \item descarte de entradas obsoletas;
    \item reapertura cuando sea necesaria;
    \item propiedades demostradas de $h$;
    \item unidades compatibles.
  \end{itemize}
  \alert{El nombre del algoritmo no sustituye estas condiciones.}
\end{frame}

% 29
\begin{frame}{Heurística admisible}
  Una heurística es admisible si nunca sobreestima:
  \[
    \boxed{0\leq h(n)\leq h^*(n)}
  \]
  \begin{itemize}
    \item Es una cota inferior, no una predicción promedio.
    \item Debe cumplirse para todos los estados relevantes.
    \item Una sola sobreestimación rompe la garantía general.
  \end{itemize}
  \begin{block}{Alcance}
    Admisibilidad se refiere al costo óptimo restante global.
  \end{block}
\end{frame}

% 30
\begin{frame}{Intuición de optimalidad}
  Para un nodo sobre un camino óptimo:
  \[
    f(n)=g(n)+h(n)\leq C^*
  \]
  Para una solución subóptima de costo $C>C^*$:
  \[
    f(G)=g(G)=C>C^*
  \]
  \begin{block}{Lectura}
    La frontera conserva antes algún candidato compatible con costo óptimo, siempre que se respeten las demás condiciones de A*.
  \end{block}
\end{frame}

% 31
\begin{frame}{Heurística consistente}
  Para toda transición $n\rightarrow n'$:
  \[
    \boxed{h(n)\leq c(n,n')+h(n')}
  \]
  Es una desigualdad triangular. Entonces:
  \[
    f(n')=g(n)+c(n,n')+h(n')\geq g(n)+h(n)=f(n)
  \]
  \begin{block}{Consecuencia}
    Con $h(G)=0$ y $h(n)\geq0$, consistencia implica la admisibilidad adoptada y los valores $f$ no disminuyen.
  \end{block}
\end{frame}

% 32
\begin{frame}{Auditoría local de consistencia}
  \small
  \begin{tabularx}{\textwidth}{>{\bfseries}p{.20\textwidth}p{.25\textwidth}>{\bfseries}p{.20\textwidth}X}
    \toprule
    Arista & Comprobación & Arista & Comprobación \\\midrule
    $S\rightarrow A$ & $7\leq2+5$ & $S\rightarrow B$ & $7\leq2+7$ \\
    $A\rightarrow C$ & $5\leq2+3$ & $A\rightarrow D$ & $5\leq5+6$ \\
    $B\rightarrow D$ & $7\leq2+6$ & $C\rightarrow G$ & $3\leq3+0$ \\
    $D\rightarrow G$ & $6\leq6+0$ & & \\
    \bottomrule
  \end{tabularx}
  \vfill
  Cero violaciones demuestra consistencia \concepto{en este grafo revisado}, no en datos futuros.
\end{frame}

% 33
\begin{frame}{Admisible pero inconsistente}
  Admisibilidad no garantiza consistencia.
  \begin{itemize}
    \item Los valores $f$ pueden disminuir a lo largo de un camino.
    \item Puede aparecer después un camino más barato hacia un estado cerrado.
    \item A* debe permitir la \concepto{reapertura} cuando mejora \texttt{mejor\_g}.
    \item Cerrar estados para siempre puede producir una solución subóptima.
  \end{itemize}
  \begin{alertblock}{Regla de implementación}
    No prohibir reaperturas sin haber demostrado consistencia.
  \end{alertblock}
\end{frame}

% 34
\begin{frame}{Cómo diseñar una heurística}
  \scriptsize
  \begin{tabularx}{\textwidth}{>{\bfseries}p{.22\textwidth}X}
    \toprule
    Principio & Idea \\\midrule
    Relajación & eliminar restricciones y resolver el problema más fácil \\
    Abstracción & buscar exactamente en un espacio reducido \\
    Geometría & usar distancias y límites físicos \\
    Descomposición & combinar subcostos sin doble conteo \\
    Patrones & precomputar costos de configuraciones abstractas \\
    \bottomrule
  \end{tabularx}
  \vfill
  \alert{La procedencia de la cota debe quedar documentada.}
\end{frame}

% 35
\begin{frame}{Combinar y evaluar heurísticas}
  \small
  Si $h_1$ y $h_2$ son admisibles:
  \begin{itemize}
    \item $\max(h_1,h_2)$ también es admisible y domina a ambas;
    \item $h_1+h_2$ requiere costos aditivos sin superposición.
  \end{itemize}
  \begin{block}{Evaluar sobre las mismas instancias}
    \begin{itemize}
      \item costo de solución;
      \item nodos generados y expandidos;
      \item frontera máxima y reaperturas;
      \item tiempo de búsqueda y costo de calcular $h$.
    \end{itemize}
  \end{block}
\end{frame}

% 36
\begin{frame}{Comparación de estrategias}
  \small
  \begin{tabularx}{\textwidth}{>{\bfseries}p{.17\textwidth}p{.18\textwidth}p{.22\textwidth}X}
    \toprule
    Estrategia & Prioridad & Óptima & Riesgo principal \\\midrule
    BFS & profundidad & costos iguales* & memoria \\
    UCS & $g$ & sí* & expansión amplia \\
    Voraz & $h$ & no & solución cara \\
    A* & $g+h$ & sí* & memoria y supuestos \\
    \bottomrule
  \end{tabularx}
  \vfill
  {\scriptsize\alert{* Bajo ramificación finita, costo mínimo positivo cuando corresponda, prueba correcta del objetivo, repetidos y propiedad heurística declarada.}}
\end{frame}

% 37
\begin{frame}{MOV-02: movilidad como grafo OD}
  \small
  Esta es una abstracción distinta de la reasignación de flota: se buscan caminos sobre conectividad OD agregada.
  \vspace{.4em}
  \begin{columns}[T]
    \column{.49\textwidth}
      \begin{block}{Representación}
        \begin{itemize}
          \item Nodo: zona TLC.
          \item Arista $i\rightarrow j$: suficientes viajes reportados.
          \item Costo: duración o distancia mediana.
          \item Referencia: UCS sobre la misma instancia.
        \end{itemize}
      \end{block}
    \column{.47\textwidth}
      \begin{block}{Heurística candidata}
        Separación geográfica con escala y unidad compatibles, auditada arista por arista.
      \end{block}
      \begin{alertblock}{Límite}
        El camino no es una ruta vial ni una trayectoria observada completa.
      \end{alertblock}
  \end{columns}
\end{frame}

% 38
\begin{frame}{Actividad: definir y probar una heurística}
  \small
  En equipos:
  \begin{enumerate}
    \item fijar origen, objetivo, costos y desempate;
    \item ejecutar UCS como referencia;
    \item proponer $h$ en la misma unidad;
    \item comprobar $h(G)=0$, no negatividad y consistencia;
    \item ejecutar A* bajo idénticas condiciones;
    \item comparar costo, camino, generados, expandidos y frontera máxima;
    \item explicar qué ocurre con $h=0$.
  \end{enumerate}
  \begin{block}{Producto}
    Definición y prueba de la heurística, sin resultados inventados.
  \end{block}
\end{frame}

% 39
\begin{frame}{Síntesis y continuidad}
  \small
  \begin{columns}[T]
    \column{.58\textwidth}
      \begin{block}{Ideas centrales}
        \begin{itemize}
          \item Voraz usa $h$; A* combina $g+h$.
          \item Admisibilidad es global; consistencia se audita localmente.
          \item A* necesita implementación y supuestos correctos.
          \item Menos expansiones no significa menor costo.
          \item $h$ comparte unidad y semántica con el costo.
        \end{itemize}
      \end{block}
    \column{.38\textwidth}
      \begin{block}{Lectura}
        Capítulo 7, §7.3\\[.4em]
        Capítulo 6, §6.4\\[.4em]
        Actividad MOV-02
      \end{block}
  \end{columns}
  \vfill
  \centering\Large\concepto{Una heurística útil también debe ser auditable.}
\end{frame}

\end{document}
